% Two ideals of Z[sqrt(-5)] drawn in the complex plane. Left: the
% principal ideal (1), the full lattice with unit cell 1 x sqrt(5).
% Right: the non-principal ideal (2, 1 + sqrt(-5)) — the sublattice of
% points a + b sqrt(-5) with a + b even — with two fundamental cells
% showing its different "shape".
% Source: book/tex/alg-NT/classgrp.tex (§ The class group) — asy figure.
\documentclass[border=2mm]{standalone}
\usepackage{tikz}
\begin{document}
\begin{tikzpicture}[scale=0.42]
  \pgfmathsetmacro{\rt}{sqrt(5)}
  % Left picture: the ideal (1).
  \begin{scope}
    \draw[dotted] (-5.2,0) -- (5.2,0);
    \draw[dotted] (0,-5.2) -- (0,5.2);
    \draw[cyan] (0,0) rectangle (1,\rt);
    \foreach \x in {-5,...,5}
      \foreach \y in {-2,...,2}
        \fill[blue] (\x,\y*\rt) circle (2.6pt);
  \end{scope}
  % Right picture: the ideal (2, 1 + sqrt(-5)).
  \begin{scope}[shift={(13.5,0)}]
    \draw[dotted] (-5.2,0) -- (5.2,0);
    \draw[dotted] (0,-5.2) -- (0,5.2);
    \draw[cyan] (0,0) rectangle (2,2*\rt);
    \draw[cyan] (0,0) -- (1,\rt) -- (-4,2*\rt) -- (-5,\rt) -- cycle;
    \foreach \x in {-5,...,5}
      \foreach \y in {-2,...,2}
      {
        \pgfmathparse{Mod(\x+\y,2)}
        \ifdim\pgfmathresult pt=0pt
          \fill[blue] (\x,\y*\rt) circle (2.6pt);
        \else
          \fill[gray!50] (\x,\y*\rt) circle (2.6pt);
        \fi
      }
  \end{scope}
\end{tikzpicture}
\end{document}
